\reading{IM-13}{Distress Symptoms and Remedies}{QUICK WIN}{2}

\srcnote{Roberto Ippolito, \textit{Private Capital Investing: The Handbook of
Private Debt and Private Equity} (Wiley, 2020), Chapter 20. No GARP source book
covers this reading.}

\begin{imtrapbox}[breakable=false,title={A reading with no chapter in either layer}]
Neither the crash course nor the lecture set contains a chapter for IM-13. The
material in this volume comes from the \term{compiled notes} supplied with the
source set, which were themselves written to the GARP learning objectives. Those
notes describe themselves as written ``from the GARP learning objectives'' --- but
since no GARP \emph{book} exists for this reading, the content sits \term{one rung
below a GARP source}. It has been checked for consistency against the available
Schweser condensation, which uses the same structure: early warning signals,
causes, and out-of-court against in-court remedies.
\end{imtrapbox}

\begin{imkeybox}[title={The three objectives in this reading}]
\begin{enumerate}[label=\textbf{\alph*.}]
\item Explain different reasons why firms may experience financial distress and identify potential indicators of distress.
\item Describe methods firms can use to remedy a financial distress situation at different stages of distress and explain the potential implications and challenges to using each method.
\item Summarize potential legal challenges that can arise when a firm is in financial distress or is going through a restructuring.
\end{enumerate}
\end{imkeybox}

% ------------------------------------------------------------------
\lo{IM-13 a}{Explain different reasons why firms may experience financial distress and identify potential indicators of distress.}

\begin{imdefbox}[title={The distinction that drives everything else}]
\term{Economic distress} means the underlying business no longer earns its cost of
capital. The assets themselves are impaired; \emph{even with zero debt the firm
would struggle}.

\smallskip
\term{Financial distress} means the business is viable but the capital structure
is not. Operating cash flow is positive; it is simply \emph{insufficient to
service the debt as structured}.
\end{imdefbox}

\begin{imkeybox}
Financial distress is fixed on the \term{right-hand side of the balance sheet},
through a restructuring. Economic distress requires an \term{operational fix} or,
failing that, liquidation --- and a debt restructuring alone only postpones the
outcome. Firms frequently suffer both at once, which is why restructurings so
often pair a debt reduction with a business plan.
\end{imkeybox}

\subsection{Indicators of distress}

\begin{itemize}
\item \term{Model-based signals.} The \term{Altman Z-score} falling through its
distress thresholds; the \term{Merton distance-to-default} shrinking; the
\term{KMV expected default frequency} rising.
\item \term{Market and behavioural indicators generally lead accounting
indicators}, because reported financials are backward-looking and published with a
lag.
\end{itemize}

\begin{imtrapbox}[breakable=false,title={Three readings converge here}]
All three model-based signals are objects this volume and its companions have
already built. The Z-score is CR-9 b; distance-to-default is the Merton model of
CR-5 e and CR-9 m; expected default frequency is the KMV model of CR-9 n. IM-13
is where they are used as \emph{symptoms} rather than as measures. And note the
lead-lag direction: \term{markets move first, accounts move later}. Reversing that
is the obvious corruption.
\end{imtrapbox}

% ------------------------------------------------------------------
\lo{IM-13 b}{Describe methods firms can use to remedy a financial distress situation at different stages of distress and explain the potential implications and challenges to using each method.}

\renewcommand{\arraystretch}{1.25}
\noindent\begin{tabularx}{\textwidth}{@{}>{\raggedright\arraybackslash}p{2.2cm}XXX@{}}
\toprule
\textbf{\sffamily\small Stage} & \textbf{\sffamily\small Remedies} & \textbf{\sffamily\small Implications} & \textbf{\sffamily\small Challenges} \\
\midrule
\textbf{1. Under\-performance} \newline \emph{liquidity still adequate}
& Operational turnaround: cost reduction, working capital release, capex deferral. Disposal of non-core assets. Suspension of dividends and buybacks. Hedging the exposure that caused the shock.
& \term{Least value destruction.} Control stays with existing management and shareholders. Nothing becomes public beyond ordinary disclosure.
& Management denial and late recognition. Cuts that repair this year can impair competitiveness for several. Asset sales in a soft market realise poor prices. \\
\textbf{3. Deep distress} \newline \emph{default likely or occurred}
& \term{Out-of-court restructuring:} distressed exchange into new debt with lower principal or longer maturity; debt-for-equity swap; consensual standstill among lenders.
& \term{Faster and far cheaper} than a court process. Confidential. Relationships with customers and suppliers largely preserved.
& \term{Holdouts.} Core payment terms usually cannot be altered without near-unanimous consent, so a small minority can block. Rating agencies treat a distressed exchange as a \term{default}. Cancellation of indebtedness generates \term{taxable income}. \\
\textbf{4. Formal insolvency}
& Prepackaged or prearranged bankruptcy. Chapter 11 reorganisation. Scheme of arrangement or restructuring plan. Receivership. Chapter 7 liquidation.
& \term{Binds dissenting creditors} through class voting and cramdown. The \term{automatic stay} halts enforcement. \term{Debtor-in-possession financing} is available with priority. Executory contracts and leases can be rejected.
& Expensive and slow. Control passes toward creditors and the court. Public disclosure. Customer and supplier relationships suffer. \\
\bottomrule
\end{tabularx}
\renewcommand{\arraystretch}{1}
\figcap{The remedy ladder, with stage 2 (liquidity stress) omitted from this
extract of the compiled notes. Read the table down the \emph{challenges} column
and the logic of the whole reading appears: the cheap remedies can be blocked by a
single holdout, and the remedy that cannot be blocked is the expensive one. That
trade-off --- \term{consent against cost} --- is the reason prepackaged
bankruptcies exist, sitting as they do in the first row of stage 4.}

\begin{imkeybox}[title={Why a distressed exchange is not a free lunch}]
Three costs attach to the cheap route. Rating agencies \term{treat it as a
default} regardless of consent. The debt forgiven generates \term{cancellation of
indebtedness income}, which is taxable. And it needs near-unanimity, so it can be
held to ransom.
\end{imkeybox}

% ------------------------------------------------------------------
\lo{IM-13 c}{Summarize potential legal challenges that can arise when a firm is in financial distress or is going through a restructuring.}

\begin{itemize}
\item \term{DIP financing and priming liens.} New money can be granted priority
\emph{over existing secured creditors} where those creditors are
``adequately protected.'' Whether protection is adequate is contested, and DIP
terms often effectively control the case.
\item \term{Ipso facto clauses.} Contract provisions that terminate on insolvency
are \term{generally unenforceable} in US bankruptcy --- though derivatives, repos
and other \term{qualified financial contracts sit within safe harbours} and can be
closed out.
\item \term{Directors in the zone of insolvency.} Fiduciary duty \term{shifts
toward creditors} as the firm approaches insolvency. Directors face wrongful or
insolvent trading exposure for continuing to incur obligations, and must document
the basis for continuing to trade.
\item \term{Cross-border recognition.} Group insolvencies raise questions of which
forum governs: recognition of foreign proceedings under the \term{UNCITRAL Model
Law} and \term{Chapter 15}, disputes over the \term{centre of main interests}, and
\term{forum shopping} toward debtor-friendly jurisdictions, all of which add cost
and delay.
\end{itemize}

\begin{imtrapbox}[title={Consolidated trap summary --- IM-13}]
\begin{itemize}
\item \term{Definition, IM-13 a.} Economic distress impairs the \emph{assets};
financial distress impairs the \emph{capital structure}. With zero debt, only the
first still has a problem.
\item \term{Polarity, IM-13 a.} Market and behavioural indicators \emph{lead};
accounting indicators \emph{lag}. Accounts are backward-looking and published late.
\item \term{Sibling, IM-13 b.} Out-of-court is cheap, fast and confidential but
blockable by holdouts. In-court binds dissenters through cramdown but is expensive,
slow and public.
\item \term{Sign, IM-13 b.} A distressed exchange is treated as a \emph{default} by
rating agencies even when every creditor consents, and the forgiven debt is
\emph{taxable income}.
\item \term{Role, IM-13 c.} Priming liens subordinate \emph{existing secured}
creditors to new DIP money, subject to adequate protection.
\item \term{Scope, IM-13 c.} Ipso facto clauses are unenforceable ---
\emph{except} for qualified financial contracts, which are safe-harboured. The
exception is the examinable half.
\item \term{Role, IM-13 c.} In the zone of insolvency, directors' duty shifts
\emph{toward creditors}, not toward shareholders.
\end{itemize}
\end{imtrapbox}

% ==================================================================
